% ============================================================================
\section{Optimización multiobjetivo de kits}
\label{sec:optimizacion}
% ============================================================================

Esta sección es el centro del documento. Con cada renglón ya resuelto a un
conjunto de candidatos, hay que elegir un candidato por renglón. Ninguna
elección es la mejor en todos los sentidos: el kit más barato no es el de
mejor marca, y el que trae el personaje favorito no es ninguno de los dos. El
problema es intrínsecamente \emph{multiobjetivo}, y el sistema lo trata como
tal. Primero calcula un \emph{frente de Pareto} ---el conjunto de kits que no
se pueden mejorar en un objetivo sin empeorar otro--- y después lo
\emph{proyecta} en los kits que se le muestran al acudiente y en uno
recomendado.

La sección está organizada así. Las \cref{sec:opt-instancia,sec:opt-unidades,sec:opt-objetivos}
formulan el problema. La \cref{sec:opt-estructura} analiza su estructura y
prueba resultados que delimitan qué se puede resolver de forma exacta y a qué
costo. Las \cref{sec:opt-exacto,sec:opt-nsga,sec:opt-escalera} describen los
tres procedimientos que construyen el frente: la enumeración exacta, NSGA-II
y la escalera de costo marginal. Las
\cref{sec:opt-proyeccion,sec:opt-lineas,sec:opt-medida,sec:opt-recomendar,sec:opt-personaje}
describen cómo el frente se convierte en kits. La \cref{sec:opt-completo}
reúne el algoritmo completo y su complejidad, y la \cref{sec:opt-evolucion}
cuenta su evolución.

% ---------------------------------------------------------------------------
\subsection{La instancia}
\label{sec:opt-instancia}
% ---------------------------------------------------------------------------

\begin{definicion}[Instancia de optimización]
\label{def:instancia}
Sea $R=(r_1,\dots,r_N)$ la secuencia de renglones del hogar, cada uno con su
hijo $h(r)$, su cantidad $q(r)$ y su conjunto de candidatos $K(r)$
(las opciones $O$ de la \cref{sec:res-niveles}, con $\card{K(r)}\leq3$). Los
candidatos \emph{admisibles} de un renglón son los que cumplen los atributos
impuestos, o todos si ninguno cumple:
\[
  A(r)=\begin{cases}\{c\in K(r):\mathrm{cumple}(c)\} & \text{si no es vacío,}\\
  K(r) & \text{en otro caso.}\end{cases}
\]
Las líneas \emph{decidibles} son
$\Dec=\{r\in R : r \text{ no es informativo},\ A(r)\neq\varnothing\}$.
Escribimos $\Dec=(r_1,\dots,r_n)$, $A_i=A(r_i)$, $k_i=\card{A_i}$ y
$k=\max_i k_i\leq3$. El \emph{espacio de búsqueda} es
\[
  \Esp=A_1\times\cdots\times A_n,\qquad \card{\Esp}=\prod_{i=1}^n k_i\leq 3^n.
\]
Una \emph{asignación} es $\mathbf{x}=(x_1,\dots,x_n)\in\Esp$.
\end{definicion}

Que un candidato que no cumple el atributo entre «como último recurso»
merece una explicación. La resolución ya se abstuvo cuando \emph{ningún}
producto del concepto cumplía. Si los candidatos llegan, es porque al menos
uno cumplía en la resolución. El caso residual (por ejemplo, un renglón
validado a mano) no deja la línea vacía: ofrece el candidato y conserva su
marca de incumplimiento visible.

% ---------------------------------------------------------------------------
\subsection{Aritmética de unidades y costo de línea}
\label{sec:opt-unidades}
% ---------------------------------------------------------------------------

\begin{definicion}[Unidades a comprar y costo de línea]
\label{def:unidades}
Para el renglón $r_i$ y el candidato $c\in A_i$,
\[
  u_i(c)=\begin{cases}\techo{\dfrac{q(r_i)}{\max(1,\emp(c))}} & q(r_i)>0,\\[4pt] 0 & q(r_i)\leq0,\end{cases}
  \qquad
  \kappa_i(c)=\begin{cases}\precio(c)\,u_i(c) & \precio(c)\neq\sinprecio,\\ \sinprecio & \text{en otro caso.}\end{cases}
\]
\end{definicion}

El $\max(1,\cdot)$ es redundante con la restricción de la base de datos
($\emp\geq1$), pero se mantiene: una división por cero en esta etapa
produciría un error en la única pantalla que el acudiente estaba esperando.

\begin{lema}[El precio unitario no ordena el costo]
\label{lem:precio-unitario}
Existen candidatos $c,c'$ y una cantidad $q$ tales que
$\precio(c)<\precio(c')$ y, sin embargo, $\kappa(c)>\kappa(c')$.
\end{lema}

\begin{proof}
Con $q=12$, $c$ un lápiz suelto ($\precio=\$1.000$, $\emp=1$) y $c'$ una caja de
doce ($\precio=\$9.000$, $\emp=12$): $\kappa(c)=\$12.000>\kappa(c')=\$9.000$.
\end{proof}

Todas las comparaciones de «más barato» del sistema se hacen por lo tanto
sobre $\kappa$ y no sobre $\precio$. El orden total que usa la heurística es
\[
  \mathrm{ord}_i(c)=\begin{cases}(0,\ \kappa_i(c),\ \precio(c),\ \mathrm{sku}(c)) & \precio(c)\neq\sinprecio,\\
  (1,\ 0,\ 0,\ \mathrm{sku}(c)) & \text{en otro caso,}\end{cases}
\]
comparado lexicográficamente: los candidatos con precio van primero, por costo
de línea, y a igual costo por precio unitario. Un candidato sin precio no se
descarta ---puede ser el único que cumple lo que el colegio exige---, pero no
puede ganar una comparación de precio.

% ---------------------------------------------------------------------------
\subsection{Objetivos y dominancia}
\label{sec:opt-objetivos}
% ---------------------------------------------------------------------------

Sea $\Pi_h$ el conjunto de personajes pedidos para el hijo $h$ (posiblemente
vacío). Un candidato \emph{satisface el personaje} de la línea $i$ si
\[
  \varphi_i(c)=\ind{\mathrm{per}(c)\neq\varnothing\ \wedge\ \bigl(\Pi_{h(r_i)}=\varnothing
  \ \vee\ \mathrm{per}(c)\in\Pi_{h(r_i)}\bigr)},
\]
comparando sin distinguir mayúsculas. Si nadie pidió un personaje, cuenta
cualquiera: el kit Fan es «el kit divertido» aunque nadie lo haya pedido
explícitamente. Si se pidió, solo cuenta el pedido.

\begin{definicion}[Objetivos]
\label{def:objetivos}
Para $\mathbf{x}\in\Esp$:
\begin{align*}
  f_1(\mathbf{x}) &= \sum_{i:\,\precio(x_i)\neq\sinprecio}\kappa_i(x_i)
    && \text{costo (a minimizar)},\\
  f_2(\mathbf{x}) &= \frac{1}{n}\sum_{i=1}^n\frac{\nivel(x_i)}{2}
    && \text{fidelidad de marca (a maximizar)},\\
  f_3(\mathbf{x}) &= \frac{1}{n}\sum_{i=1}^n\ind{\precio(x_i)\neq\sinprecio}
    && \text{disponibilidad (a maximizar)},\\
  f_4(\mathbf{x}) &= \frac{1}{n}\sum_{i=1}^n\varphi_i(x_i)
    && \text{personaje (a maximizar)}.
\end{align*}
Los tres últimos se redondean a seis decimales para que la comparación entre
asignaciones sea exacta y reproducible.
\end{definicion}

Hay tres decisiones de modelado en estas definiciones:
\begin{itemize}
  \item \textbf{Los candidatos sin precio no entran al costo.} Su ausencia se
        refleja en $f_3$. Contarlos como cero abarataría el kit de forma
        engañosa; contarlos como infinito los haría inelegibles. La
        \cref{obs:sin-precio} discute una consecuencia de esta decisión.
  \item \textbf{La fidelidad es el nivel de marca promedio normalizado a
        $[0,1]$.} Es lo que separa un kit Premium de uno Ahorrador.
  \item \textbf{Las líneas informativas no participan.} Medirlas inflaría los
        promedios con ítems que el sistema nunca iba a surtir, y el kit se
        vería mejor cuanto menos pudiera vender.
\end{itemize}

\begin{definicion}[Dominancia y frente de Pareto]
\label{def:dominancia}
Sea $\mathbf{F}(\mathbf{x})=(f_1,f_2,f_3,f_4)(\mathbf{x})$. Decimos que
$\mathbf{x}$ \emph{domina} a $\mathbf{y}$, $\mathbf{x}\domina\mathbf{y}$, si
\[
  f_1(\mathbf{x})\leq f_1(\mathbf{y}),\ \ f_j(\mathbf{x})\geq f_j(\mathbf{y})\ (j=2,3,4),
  \ \ \text{y al menos una desigualdad es estricta.}
\]
El \emph{conjunto de Pareto} es $\Esp^\ast=\{\mathbf{x}\in\Esp:\nexists\,\mathbf{y}\in\Esp,\ \mathbf{y}\domina\mathbf{x}\}$
y el \emph{frente de Pareto} es su imagen $\Front=\mathbf{F}(\Esp^\ast)$
\cite{miettinen1999,ehrgott2005}. En la implementación, el frente se representa
con una asignación por cada vector objetivo distinto, ordenado por
$(f_1,-f_2,-f_3,-f_4)$ para que su presentación sea estable.
\end{definicion}

% ---------------------------------------------------------------------------
\subsection{Estructura del problema}
\label{sec:opt-estructura}
% ---------------------------------------------------------------------------

Antes de elegir un algoritmo conviene entender qué tan difícil es el
problema. El análisis que sigue muestra que la formulación vigente tiene una
estructura muy favorable y explica los resultados empíricos de la
\cref{sec:evaluacion}.

\begin{proposicion}[Separabilidad aditiva]
\label{prop:separable}
Cada objetivo es una suma de términos que dependen de una sola línea:
$f_j(\mathbf{x})=\sum_{i=1}^n g_{j,i}(x_i)$ para funciones
$g_{j,i}\colon A_i\to\R$ con
$g_{1,i}(c)=\kappa_i(c)\ind{\precio(c)\neq\sinprecio}$,
$g_{2,i}(c)=\nivel(c)/(2n)$, $g_{3,i}(c)=\ind{\precio(c)\neq\sinprecio}/n$ y
$g_{4,i}(c)=\varphi_i(c)/n$.
\end{proposicion}

\begin{proof}
Es inmediato de la \cref{def:objetivos}. Ningún término depende de la
elección hecha en otra línea, en particular porque la consolidación entre
hermanos (\cref{sec:opt-lineas}) se aplica \emph{después} de elegir y no forma
parte de $f_1$.
\end{proof}

\begin{corolario}[Los extremos se resuelven línea por línea]
\label{cor:extremos}
Para cualquier orden lexicográfico de los objetivos, por ejemplo
\emph{mínimo costo y, a igual costo, máxima disponibilidad}, la asignación
óptima se obtiene eligiendo en cada línea, de forma independiente, el
candidato óptimo bajo ese orden. El costo es $O(\sum_i k_i)=O(nk)$.
\end{corolario}

\begin{proof}
Por la \cref{prop:separable}, $\min f_1=\sum_i\min_{c\in A_i}g_{1,i}(c)$, y
el conjunto de minimizadores es el producto cartesiano de los minimizadores
por línea. Dentro de ese producto, el segundo criterio también es separable,
y el argumento se repite.
\end{proof}

Esto significa que los kits de los extremos ---el más barato, el de mejor
marca y el de más personaje--- no requieren ninguna búsqueda: un algoritmo
voraz que mira cada línea por separado es óptimo. El problema solo se vuelve
interesante en los \emph{compromisos}, por ejemplo la mejor marca que cabe en
un presupuesto.

\begin{proposicion}[El compromiso con presupuesto es una mochila de elección múltiple]
\label{prop:mckp}
El problema
\[
  \max_{\mathbf{x}\in\Esp} f_2(\mathbf{x})\quad\text{sujeto a}\quad f_1(\mathbf{x})\leq B
\]
es una instancia del problema de la mochila de elección múltiple (MCKP)
\cite{sinha1979mckp,kellerer2004}: hay $n$ clases disjuntas $A_i$, de cada
clase se elige exactamente un ítem con peso $g_{1,i}$ y valor
$g_{2,i}$, y se maximiza el valor total con el peso acotado por $B$.
\end{proposicion}

En general, el MCKP es NP-difícil: contiene a la mochila 0-1, cuya versión de
decisión es NP-completa \cite{garey1979}, como el caso particular de clases
$\{\text{tomar}, \text{no tomar}\}$. Sin embargo, el MCKP es NP-difícil solo en
sentido \emph{débil}: admite algoritmos pseudopolinomiales cuyo costo depende
de la magnitud de los valores o de los pesos \cite{kellerer2004}. En esta
instancia, los valores tienen una magnitud muy pequeña:
\[
  n\,f_2(\mathbf{x})\cdot 2=\sum_i\nivel(x_i)\in\{0,1,\dots,2n\}.
\]

\begin{teorema}[El compromiso con presupuesto es polinomial]
\label{teo:mckp-poli}
El problema de la \cref{prop:mckp} se resuelve de forma exacta en tiempo
$O(n^2k)$ y espacio $O(n)$ por estado, mediante la recurrencia
\[
  M_0(0)=0,\qquad M_0(s)=+\infty\ (s\neq0),\qquad
  M_i(s)=\min_{c\in A_i,\ \nivel(c)\leq s}\Bigl\{M_{i-1}\bigl(s-\nivel(c)\bigr)+g_{1,i}(c)\Bigr\},
\]
donde $M_i(s)$ es el mínimo costo de una asignación parcial de las primeras
$i$ líneas cuya suma de niveles es exactamente $s$. El óptimo es
$s^\ast=\max\{s : M_n(s)\leq B\}$, con valor $f_2^\ast=s^\ast/(2n)$, y la
asignación se recupera guardando el argumento que alcanza cada mínimo.
\end{teorema}

\begin{proof}
\emph{Corrección}: por inducción sobre $i$. Toda asignación parcial de suma
$s$ termina en algún $c\in A_i$ y su prefijo tiene suma $s-\nivel(c)$, con
costo al menos $M_{i-1}(s-\nivel(c))$ por hipótesis inductiva. Esa cota se
alcanza extendiendo el prefijo óptimo. Como $f_2$ es una función creciente de
$s$, maximizar $f_2$ con $f_1\leq B$ equivale a encontrar el mayor $s$
alcanzable con costo mínimo a lo sumo $B$. \emph{Costo}: hay $n$ etapas,
cada una con a lo sumo $2n+1$ estados y $k$ transiciones por estado.
\end{proof}

Esta recurrencia es exactamente la que se usó como verdad de referencia en la
evaluación (\cref{sec:evaluacion}). El mismo argumento se extiende al frente
completo de los cuatro objetivos, porque los tres objetivos «de conteo»
también toman pocos valores.

\begin{teorema}[El frente exacto se calcula en tiempo polinomial]
\label{teo:frente-poli}
Sean $a(\mathbf{x})=\sum_i\nivel(x_i)$, $b(\mathbf{x})=\sum_i\ind{\precio(x_i)\neq\sinprecio}$
y $p(\mathbf{x})=\sum_i\varphi_i(x_i)$, de modo que
$(f_2,f_3,f_4)=\bigl(a/(2n),\ b/n,\ p/n\bigr)$. Sea
\[
  C^\ast(a,b,p)=\min\{f_1(\mathbf{x}) : \mathbf{x}\in\Esp,\ (a,b,p)(\mathbf{x})=(a,b,p)\},
\]
calculado por la recurrencia de la \cref{teo:mckp-poli} sobre estados
$(a,b,p)\in[0,2n]\times[0,n]\times[0,n]$. Entonces:
\begin{enumerate}
  \item $\card{\Front}\leq(2n+1)(n+1)^2$.
  \item $\Front$ es el conjunto de vectores no dominados de
        $\bigl\{\bigl(C^\ast(a,b,p),\tfrac{a}{2n},\tfrac bn,\tfrac pn\bigr)\bigr\}$.
  \item El frente se calcula en tiempo $O(n^4k)$.
\end{enumerate}
\end{teorema}

\begin{proof}
(1) Dos puntos distintos del frente no pueden compartir $(a,b,p)$: si lo
hicieran, el de mayor costo estaría dominado por el otro. El número de ternas
posibles es $(2n+1)(n+1)^2$.
(2) Sea $\mathbf{x}\in\Esp^\ast$ con terna $(a,b,p)$. Por definición,
$C^\ast(a,b,p)\leq f_1(\mathbf{x})$. Si la desigualdad fuera estricta, la
asignación que alcanza $C^\ast$ tendría los mismos $f_2,f_3,f_4$ y menor
costo, y dominaría a $\mathbf{x}$, lo que es una contradicción. Por lo tanto,
todo punto del frente está en el conjunto de los $C^\ast$. Recíprocamente, sea
$\mathbf{v}$ un punto de ese conjunto que no está dominado dentro de él, y
supóngase que algún $\mathbf{y}\in\Esp$ lo domina. El punto $C^\ast$ de la
terna de $\mathbf{y}$ tiene los mismos $f_2,f_3,f_4$ que $\mathbf{y}$ y un
costo no mayor, así que también domina a $\mathbf{v}$, lo que contradice la
hipótesis.
(3) Tras procesar $i$ líneas hay a lo sumo $(2i+1)(i+1)^2=O(i^3)$ estados,
cada uno con $k$ transiciones, así que la recurrencia cuesta
$O\bigl(k\sum_{i\leq n}i^3\bigr)=O(n^4k)$. Para el filtro final, el punto de
terna $t$ está dominado si y solo si existe $t'\geq t$ componente a componente,
con $t'\neq t$ y $C^\ast(t')\leq C^\ast(t)$. El mínimo de $C^\ast$ sobre ese
ortante se obtiene para todas las ternas con el mínimo acumulado
$\mathrm{Suf}(t)=\min\{C^\ast(t),\mathrm{Suf}(t+e_1),\mathrm{Suf}(t+e_2),\mathrm{Suf}(t+e_3)\}$,
en $O(n^3)$, así que el filtro no cambia el orden.
\end{proof}

\begin{observacion}[Verificación empírica]
Se comparó una implementación directa de esta programación dinámica con la
enumeración exacta del sistema (\cref{sec:opt-exacto}) en 300 instancias
aleatorias de 2 a 6 líneas, con candidatos de nivel, precio (incluido
$\sinprecio$), empaque y personaje aleatorios. Los dos frentes coincidieron
en las 300 instancias. Para una lista típica ($n=17$, $k=3$), el número de
estados es como máximo $35\cdot18^2=\num{11340}$, frente a
$\card{\Esp}=3^{17}\approx1{,}3\times10^8$.
\end{observacion}

\begin{hallazgo}{Consecuencia para el diseño del optimizador}
Bajo la formulación vigente, el frente de Pareto exacto se puede calcular en
tiempo polinomial y con pocos estados. El NSGA-II
(\cref{sec:opt-nsga}) es por lo tanto un método \emph{aproximado} aplicado a
un problema que admite solución exacta. Eso explica dos resultados de la
evaluación. Primero, en los extremos el NSGA-II iguala siempre al óptimo,
pero una regla voraz corregida también lo hace (\cref{cor:extremos}).
Segundo, el frente del NSGA-II no es exacto: una regla simple dominaba algún
punto suyo en el 68\,\% de las listas antes de las correcciones de la
\cref{sec:opt-escalera}.

La justificación original del NSGA-II era un \emph{acoplamiento} entre líneas
que la formulación vigente no tiene. El diseño proponía, entre otras cosas,
un costo fijo por tienda utilizada, que depende de la asignación completa, e
incorporar la consolidación entre hermanos \emph{dentro} de $f_1$, de modo que
el costo de un producto dependa de la demanda agregada de todo el hogar. Con
cualquiera de esos dos términos, la \cref{prop:separable} deja de valer y el
problema deja de descomponerse. Mientras no se incorporen, la programación
dinámica del \cref{teo:frente-poli} es una alternativa exacta y más barata
(\cref{sec:discusion}).
\end{hallazgo}

% ---------------------------------------------------------------------------
\subsection{Enumeración exacta para espacios pequeños}
\label{sec:opt-exacto}
% ---------------------------------------------------------------------------

Cuando $\card{\Esp}\leq\Lambda_{\mathrm{enum}}=4096$, el sistema enumera el
espacio completo y filtra los no dominados (\cref{alg:frente-exacto}). Se
eligió una potencia de dos con un criterio de costo: $4096=2^{12}$
asignaciones se evalúan y se filtran en milisegundos. Si $\card{\Esp}\leq1$
---ninguna línea decidible, o todas con un único candidato---, no hay
compromiso que optimizar y el frente se declara vacío. En ese caso, los kits
salen de la heurística por perfil, que en ese escenario coincide con el
óptimo.

\begin{algorithm}[htbp]
\caption{Frente exacto por enumeración y filtro de no dominancia}
\label{alg:frente-exacto}
\begin{algorithmic}[1]
\Require líneas decidibles $\Dec$; admisibles $A_1,\dots,A_n$
\Ensure frente $\Front$ (una asignación por vector objetivo)
\State $S\gets\{(\mathbf{F}(\mathbf{x}),\mathbf{x}) : \mathbf{x}\in A_1\times\cdots\times A_n\}$
\State $U\gets$ un representante por vector objetivo distinto de $S$ \Comment{deduplicación}
\State $\Front\gets\{(\mathbf{v},\mathbf{x})\in U : \nexists(\mathbf{w},\cdot)\in U,\ \mathbf{w}\domina\mathbf{v}\}$
\State ordenar $\Front$ por $(f_1,-f_2,-f_3,-f_4)$
\State \Return $\Front$
\end{algorithmic}
\end{algorithm}

\paragraph{Complejidad.} Evaluar las $\card{\Esp}$ asignaciones cuesta
$O(\card{\Esp}\,n)$. El filtro por pares cuesta $O(\card{U}^2M)$ con $M=4$
objetivos. En el peor caso admitido son unos $1{,}7\times10^7$ comparaciones,
que se resuelven en un tiempo menor que el de una sola llamada de red. La
corrección es inmediata: el algoritmo aplica literalmente la
\cref{def:dominancia} sobre todo el espacio.

% ---------------------------------------------------------------------------
\subsection{NSGA-II para espacios grandes}
\label{sec:opt-nsga}
% ---------------------------------------------------------------------------

Cuando $\card{\Esp}>4096$ ---en el corpus, 137 de cada 175 listas---, el frente
se aproxima con NSGA-II \cite{deb2002nsga2}, en la implementación de la
biblioteca \emph{pymoo} \cite{blank2020pymoo}. Esta subsección describe el
algoritmo en general y después cada decisión específica de su uso aquí.

\subsubsection{El algoritmo}

NSGA-II es un algoritmo evolutivo elitista que mantiene una población $P_t$ de
$N$ soluciones y, en cada generación, produce $N$ descendientes $Q_t$ por
selección, cruce y mutación. La población siguiente se elige de
$R_t=P_t\cup Q_t$ con dos criterios: el \emph{rango de no dominancia} y la
\emph{distancia de aglomeración}.

\begin{definicion}[Ordenamiento no dominado]
Los \emph{frentes} $\mathcal{F}_1,\mathcal{F}_2,\dots$ de un conjunto $R$ se
definen recursivamente: $\mathcal{F}_1$ son los no dominados de $R$ y
$\mathcal{F}_{j+1}$ son los no dominados de
$R\setminus(\mathcal{F}_1\cup\cdots\cup\mathcal{F}_j)$. El \emph{rango} de una
solución es el índice de su frente. El algoritmo rápido de Deb et al.
\cite{deb2002nsga2} calcula todos los frentes en $O(M\card{R}^2)$, contando
para cada solución cuántas la dominan y a cuáles domina.
\end{definicion}

\begin{definicion}[Distancia de aglomeración]
Dentro de un frente $\mathcal{F}$, se ordenan las soluciones por cada objetivo
$m$. A las dos soluciones extremas se les asigna distancia $+\infty$, y a cada
solución interior $i$ se le suma
\[
  \mathrm{cd}(i)\mathrel{+}=\frac{f_m^{(i+1)}-f_m^{(i-1)}}{f_m^{\max}-f_m^{\min}},
\]
donde $(i-1)$ e $(i+1)$ son sus vecinas en el orden del objetivo $m$. La
distancia mide qué tan aislada está una solución en el espacio de objetivos, y
se calcula en $O(M\card{\mathcal{F}}\log\card{\mathcal{F}})$.
\end{definicion}

El \emph{orden de aglomeración} prefiere el menor rango y, a igual rango, la
mayor distancia. La selección de padres es por \emph{torneo binario} con ese
orden, y la supervivencia toma frentes completos en orden de rango hasta
llenar $N$, desempatando el último frente por distancia. El resultado es la
aproximación $\mathcal{F}_1$ de la última generación.

\begin{algorithm}[htbp]
\caption{NSGA-II (esquema general)}
\label{alg:nsga2}
\begin{algorithmic}[1]
\Require tamaño $N$; generaciones $G$; población inicial $P_0$; semilla $\varsigma$
\Ensure aproximación del frente
\State evaluar $\mathbf{F}$ en $P_0$
\For{$t\gets0$ \textbf{hasta} $G-1$}
  \State $Q_t\gets\emptyset$
  \While{$\card{Q_t}<N$}
    \State elegir dos padres por torneo binario con el orden de aglomeración
    \State producir dos hijos por cruce SBX y mutación polinomial; redondear y reparar
    \State $Q_t\gets Q_t\cup\{\text{hijos que no estén duplicados}\}$
  \EndWhile
  \State evaluar $\mathbf{F}$ en $Q_t$;\quad $R_t\gets P_t\cup Q_t$
  \State $(\mathcal{F}_1,\mathcal{F}_2,\dots)\gets$ ordenamiento no dominado de $R_t$
  \State $P_{t+1}\gets\emptyset$;\quad $j\gets1$
  \While{$\card{P_{t+1}}+\card{\mathcal{F}_j}\leq N$}
    \State calcular la distancia de aglomeración en $\mathcal{F}_j$;\quad $P_{t+1}\gets P_{t+1}\cup\mathcal{F}_j$;\quad $j\gets j+1$
  \EndWhile
  \State completar $P_{t+1}$ con los elementos de $\mathcal{F}_j$ de mayor distancia
\EndFor
\State \Return los no dominados de $P_G$
\end{algorithmic}
\end{algorithm}

\subsubsection{Codificación}

La variable $x_i$ es directamente el \emph{índice} del candidato elegido en la
línea $i$: $x_i\in\{0,\dots,k_i-1\}$, con cotas $x_i^{\mathrm{L}}=0$ y
$x_i^{\mathrm{U}}=k_i-1$. Esta codificación tiene una propiedad importante:
\emph{todo} vector entero dentro de las cotas es un kit válido. No hay
restricciones que reparar y el problema se declara sin restricciones.

pymoo minimiza, así que el vector que se optimiza es
$(f_1,-f_2,-f_3,-f_4)$. La evaluación reutiliza literalmente la misma función
de objetivos que la enumeración exacta y que la heurística. Tener una sola
definición de los objetivos es una decisión deliberada: una copia vectorizada
correría el riesgo de divergir de ella, que es precisamente la clase de
defecto que hace que un frente «óptimo» no sea comparable con nada.

\subsubsection{Operadores}

\paragraph{Cruce binario simulado (SBX).} Con probabilidad $p_c=0{,}9$ por
pareja, para cada variable se toma $u\sim\mathcal{U}(0,1)$ y
\[
  \beta_q=\begin{cases}(2u)^{1/(\eta_c+1)} & u\leq\tfrac12,\\[2pt]
  \bigl(\tfrac{1}{2(1-u)}\bigr)^{1/(\eta_c+1)} & u>\tfrac12,\end{cases}
  \qquad
  y^{(1,2)}=\tfrac12\bigl[(x^{(1)}+x^{(2)})\mp\beta_q\,\abs{x^{(2)}-x^{(1)}}\bigr],
\]
con índice de distribución $\eta_c=15$ \cite{deb1995sbx}. Un $\eta$ grande
concentra a los hijos cerca de los padres.

\paragraph{Mutación polinomial (PM).} Cada individuo se muta con probabilidad
$1$ y, dentro de él, cada variable con probabilidad $1/n$ (el valor por
omisión de la biblioteca). Para una variable,
\[
  \delta_q=\begin{cases}(2u)^{1/(\eta_m+1)}-1 & u<\tfrac12,\\
  1-\bigl(2(1-u)\bigr)^{1/(\eta_m+1)} & u\geq\tfrac12,\end{cases}
  \qquad y=x+\delta_q\,(x^{\mathrm{U}}-x^{\mathrm{L}}),
\]
con $\eta_m=20$ \cite{deb1996pm}. La biblioteca usa la variante acotada, que
ajusta $\delta_q$ según la distancia a cada cota.

\paragraph{Reparación por redondeo.} Los dos operadores producen valores
reales. Después de cruce y mutación, cada variable se redondea al entero más
cercano y se recorta a $[x^{\mathrm{L}},x^{\mathrm{U}}]$. La evaluación vuelve
a recortar por precaución, para no depender de esa invariante externa. Se
eliminan además los individuos duplicados.

\begin{hallazgo}{Sesgo ordinal sobre una variable nominal}
SBX y PM son operadores \emph{ordinales}: suponen que índices vecinos
representan soluciones vecinas, y en particular que el índice 3 está «entre»
el 2 y el 4. Aquí el índice es una \emph{etiqueta nominal}, y el orden de los
candidatos dentro de una línea no tiene significado en el espacio de
objetivos. Los operadores siguen funcionando como perturbaciones, pero su
sesgo hacia valores cercanos no corresponde a ninguna estructura real del
problema. La evaluación midió el efecto: antes de las correcciones, una regla
de subida por costo marginal dominaba algún punto del frente en el 68\,\% de
las listas. El sesgo se corrigió \emph{sembrando} la población
(\cref{sec:opt-siembra}), sin cambiar los operadores. Un cruce uniforme por
variable y una mutación que reemplace el índice por otro al azar serían los
operadores naturales para esta codificación.
\end{hallazgo}

\subsubsection{Tamaño de población y número de generaciones}

Los dos parámetros se escalan con el tamaño del problema, para que el frente no
se quede corto en una lista larga ni se gasten generaciones de más en una
corta:
\begin{align*}
  N(n)&=N'+(N'\bmod 2),\qquad N'=\min\bigl(200,\ \max(40,\ 12n)\bigr),\\
  G(\Esp)&=\piso{\min\Bigl(120,\ \max\bigl(30,\ 20\log_{10}\max(10,\card{\Esp})\bigr)\Bigr)}.
\end{align*}
$N$ se redondea a par porque los padres se emparejan de a dos y una población
impar desperdicia un individuo por generación. $G$ crece con el
\emph{logaritmo} del espacio porque el espacio crece multiplicativamente
---es un producto de $k_i$--- y el esfuerzo útil crece con el número de
dígitos, no con el tamaño. Para la lista típica ($n=17$, $\bar k=2{,}4$,
$\card{\Esp}\approx2{,}9\times10^6$) se obtiene $N=200$ y $G=120$, es decir,
los dos topes.

\subsubsection{Siembra de la población inicial}
\label{sec:opt-siembra}

La población inicial no es aleatoria. Se siembra con:
\begin{enumerate}
  \item las asignaciones de la heurística por perfil
        (\cref{def:heuristica}) para cada perfil, que garantizan los extremos
        de mínimo costo, máximo nivel y máximo personaje; y
  \item los peldaños de la escalera de costo marginal
        (\cref{sec:opt-escalera}), que recorren el compromiso entre costo y
        fidelidad entre esos extremos.
\end{enumerate}
Las semillas se deduplican. Si hay más semillas que cupos, se toman $N$ de
ellas espaciadas uniformemente a lo largo de la lista (índices
$\mathrm{round}(\mathrm{linspace}(0,\card{S}-1,N))$), de modo que se conserva
el recorrido completo de barato a caro. El resto de la población se completa
con asignaciones aleatorias generadas con un generador inicializado con la
misma semilla del algoritmo.

\begin{proposicion}[El frente no queda por debajo de sus semillas]
\label{prop:siembra}
Sea $\Front_{\mathrm{G}}$ el conjunto devuelto por el NSGA-II y $S$ el
conjunto de semillas. El sistema devuelve
$\Front=\mathrm{NoDom}(\Front_{\mathrm{G}}\cup S)$. Entonces ninguna semilla
domina a un punto de $\Front$, y toda semilla está en $\Front$ o está dominada
por un punto de $\Front$.
\end{proposicion}

\begin{proof}
Toda semilla pertenece al conjunto sobre el que se aplica el filtro. Si una
semilla dominara a un punto $\mathbf{v}$, ese punto habría sido eliminado y no
estaría en $\Front$. Si una semilla no está en $\Front$, es porque algún punto
del conjunto la domina. Como la dominancia es transitiva y el conjunto es
finito, una cadena de dominancia termina en un no dominado, que está en
$\Front$.
\end{proof}

La consecuencia práctica es una garantía de \emph{no regresión}: el NSGA-II
puede perder una semilla buena durante la evolución, pero el resultado final
nunca es peor que la heurística ni que la escalera. En particular, el frente
nunca es dominado por la heurística, que es la referencia mínima que se le
exige.

\subsubsection{Reproducibilidad}

La terminación es por número fijo de generaciones $G$, no por un criterio de
convergencia. Un criterio de convergencia hace depender el punto de corte del
azar de la corrida, y el requisito es que una semilla dada reproduzca el
frente byte a byte. Por la misma razón, la población aleatoria inicial usa la
semilla del algoritmo, y el frente final se ordena de forma total. Empíricamente
(\cref{sec:evaluacion}), con cinco semillas distintas el kit resultante es el
mismo en el 94,5\,\% de las listas donde corre el NSGA-II, y siempre en las que
se enumeran.

\subsubsection{Complejidad}

Por generación, el ordenamiento no dominado cuesta $O(M(2N)^2)$; la
distancia de aglomeración, $O(M\,2N\log 2N)$; y la evaluación de los $N$
hijos, $O(Nn)$. El total es
\[
  T_{\mathrm{NSGA}}=O\bigl(G\,(MN^2+Nn)\bigr),
\]
que con $G=120$, $N=200$, $M=4$ y $n=17$ da del orden de $10^7$ operaciones
elementales. El tiempo medido para producir los kits de una lista es de unos
5~s (\cref{sec:evaluacion}). Compárese con los $O(n^4k)$ estados de la
programación dinámica exacta (\cref{teo:frente-poli}), que para el mismo
tamaño son unos $10^4$ estados.

% ---------------------------------------------------------------------------
\subsection{La escalera de costo marginal}
\label{sec:opt-escalera}
% ---------------------------------------------------------------------------

La escalera es un procedimiento voraz que recorre el compromiso entre costo y
fidelidad. Parte del kit más barato y, en cada paso, sube \emph{una} línea a un
candidato de mejor nivel: la que tenga el menor costo extra por nivel ganado.

\begin{algorithm}[htbp]
\caption{Escalera de costo marginal}
\label{alg:escalera}
\begin{algorithmic}[1]
\Require líneas $\Dec$; admisibles $A_1,\dots,A_n$
\Ensure peldaños $\mathbf{x}^{(0)},\mathbf{x}^{(1)},\dots$
\For{$i\gets1$ \textbf{hasta} $n$}
  \State $x_i\gets\argmin_{c\in A_i}\bigl(\mathrm{ord}_i(c),\ -\nivel(c),\ \mathrm{sku}(c)\bigr)$ \Comment{kit más barato}
\EndFor
\State $\mathcal{S}\gets(\mathbf{x})$
\Loop
  \State $\mathit{mejor}\gets\varnothing$
  \ForAll{$i$ y $c\in A_i$ con $\precio(c)\neq\sinprecio$ y $\nivel(c)>\nivel(x_i)$}
    \State $\Delta\gets\kappa_i(c)-\kappa_i(x_i)$ \Comment{$\kappa_i(x_i)$ se toma como $0$ si es $\sinprecio$}
    \State $e\gets\Delta/(\nivel(c)-\nivel(x_i))$ \Comment{costo por nivel ganado}
    \State $\mathit{mejor}\gets\min_{\mathrm{lex}}\bigl(\mathit{mejor},\ (e,\Delta,i,c)\bigr)$
  \EndFor
  \If{$\mathit{mejor}=\varnothing$} \Return $\mathcal{S}$ \EndIf
  \State $(\cdot,\cdot,i,c)\gets\mathit{mejor}$;\quad $x_i\gets c$;\quad $\mathcal{S}\gets\mathcal{S}\cdot(\mathbf{x})$
\EndLoop
\end{algorithmic}
\end{algorithm}

\begin{proposicion}[Terminación y costo]
\label{prop:escalera}
\Cref{alg:escalera} produce como máximo $2n+1$ peldaños y corre en
$O(n^2k)$.
\end{proposicion}

\begin{proof}
Cada paso aumenta estrictamente $a(\mathbf{x})=\sum_i\nivel(x_i)$, porque solo
se aceptan candidatos de nivel mayor. Como $0\leq a\leq2n$, hay a lo sumo $2n$
pasos. Cada paso examina $O(nk)$ pares.
\end{proof}

El criterio $e=\Delta/\Delta\nivel$ es el de \emph{eficiencia incremental} del
algoritmo voraz clásico para la relajación lineal del MCKP
\cite{sinha1979mckp}. Sin embargo, la escalera no implementa ese algoritmo de
forma exacta: no descarta previamente los candidatos dominados por la
envolvente convexa (los «LP-dominados»), así que sus peldaños aproximan la
envolvente convexa inferior de la curva costo-fidelidad sin garantizar que
estén sobre ella. Su papel no es ser exacta sino complementar al NSGA-II:
recorre precisamente la región que los operadores ordinales no encuentran
bien, y por eso siembra la población y entra al filtro final. Con esa
siembra, en la evaluación la escalera dejó de dominar puntos del frente
(de 68\,\% a 0\,\%).

% ---------------------------------------------------------------------------
\subsection{La heurística por perfil y la proyección del frente}
\label{sec:opt-proyeccion}
% ---------------------------------------------------------------------------

\begin{definicion}[Heurística por perfil]
\label{def:heuristica}
Para una línea con admisibles $A_i$ y cantidad $q_i$:
\begin{itemize}
  \item \textbf{Ahorrador}: $\argmin_{c\in A_i}(\mathrm{ord}_i(c),\mathrm{sku}(c))$.
  \item \textbf{Premium}: entre los de nivel máximo $\max_{c\in A_i}\nivel(c)$,
        el de menor $(\mathrm{ord}_i,\mathrm{sku})$.
  \item \textbf{Fan}: entre los que satisfacen el personaje ($\varphi_i=1$), el
        de menor $(\mathrm{ord}_i,\mathrm{sku})$; si no hay ninguno, el
        Ahorrador. La heurística no inventa un personaje.
\end{itemize}
\end{definicion}

La heurística tiene hoy dos usos acotados: elegir en las líneas sin
compromiso, cuando el frente es vacío, y servir de semilla y de referencia
mínima para el frente. No es un plan alternativo para cuando falte la
biblioteca de optimización. Esa biblioteca es una dependencia obligatoria: si
falta, el servicio no arranca, porque una degradación silenciosa a la
heurística serviría kits distintos de los del algoritmo sin avisar a nadie.

\begin{definicion}[Proyección del frente en perfiles]
\label{def:proyeccion}
Dado un frente no vacío $\Front$, sea
\[
  \Front_{\mathrm{s}}=\Bigl\{\mathbf{x}\in\Front : f_3(\mathbf{x})=\max_{\mathbf{y}\in\Front}f_3(\mathbf{y})\Bigr\}
\]
el subconjunto de puntos \emph{surtibles}, los de máxima disponibilidad. Entonces
\begin{align*}
  \mathbf{x}_{\mathrm{Ah}} &= \lexmin_{\mathbf{x}\in\Front_{\mathrm{s}}}\ \bigl(f_1,\ -f_3,\ -f_2\bigr),\\
  \mathbf{x}_{\mathrm{Pr}} &= \lexmax_{\mathbf{x}\in\Front_{\mathrm{s}}}\ \bigl(f_2,\ -f_1,\ f_3\bigr),\\
  \mathbf{x}_{\mathrm{Fan}} &= \begin{cases}\lexmax_{\mathbf{x}\in\Front_{\mathrm{s}}}(f_4,-f_1,f_2) & \text{si ese máximo tiene } f_4>0,\\
  \mathbf{x}_{\mathrm{Ah}} & \text{en otro caso.}\end{cases}
\end{align*}
\end{definicion}

Los perfiles son \emph{arquetipos}: cada uno proyecta el frente sobre un eje
distinto. Los pesos que declara el acudiente no intervienen aquí; solo deciden
cuál de los kits se recomienda (\cref{sec:opt-recomendar}). Como los tres kits
son puntos del mismo frente, son comparables entre sí con el mismo vocabulario
de objetivos.

\paragraph{Por qué se proyecta sobre los puntos surtibles.}
\label{obs:sin-precio}
Como $f_1$ excluye las líneas sin precio, cambiar un candidato con precio por
uno sin precio \emph{reduce} $f_1$ y reduce $f_3$. Ninguno de los dos puntos
domina al otro, así que ambos pueden estar en el frente. Una versión anterior
proyectaba sobre todo $\Front$, y entonces un candidato sin precio parecía
gratis. Con dos líneas, $A_1=\{\text{económico a }\$3.000,\ \text{premium sin
precio}\}$ y $A_2=\{\text{económico a }\$2.000,\ \text{premium a }\$9.000\}$,
el Ahorrador elegía el premium sin precio y reportaba un total de \$2.000 que
no incluía esa línea. El Premium y A~tu~Medida tenían el mismo problema: una
marca sin precio sumaba fidelidad sin sumar costo. La heurística nunca lo tuvo,
porque pone lo que no tiene precio al final de su orden ($\mathrm{ord}_i$).
Restringir la proyección a $\Front_{\mathrm{s}}$ aplica esa misma regla al
frente. Lo que no se puede surtir no gana una comparación de precio, pero si
una línea \emph{solo} tiene candidatos sin precio, todos los puntos los llevan
y la línea sigue apareciendo en el kit. El frente en sí no cambia: sigue
siendo el de cuatro objetivos, y el filtro forma parte de la proyección. En el
catálogo vigente todos los productos activos tienen precio, así que la
corrección no cambia ningún kit actual; evita que el defecto aparezca cuando
una oferta venza.

\begin{proposicion}[Un kit proyectado no deja de surtir lo que podía surtir]
\label{prop:surtibles}
En toda línea cuyos admisibles incluyen al menos un candidato con precio, los
kits Ahorrador, Premium, Fan y A~tu~Medida proyectados eligen un candidato
con precio.
\end{proposicion}

\begin{proof}
Sea $b^\ast$ el número de líneas que admiten un candidato con precio; ninguna
asignación tiene más de $b^\ast$ líneas con precio. Un óptimo lexicográfico de
$(f_3,-f_1,f_2,f_4)$ es eficiente: cualquier asignación que lo dominara sería
al menos igual en $f_3$, y mejor o igual en el resto, contradiciendo la
optimalidad lexicográfica. Ese óptimo tiene $f_3=b^\ast/n$, así que hay un
punto del frente exacto con esa disponibilidad. En el NSGA-II, la semilla
Ahorrador de la heurística también tiene $f_3=b^\ast/n$, porque pone los
candidatos con precio primero, y entra al filtro final
(\cref{prop:siembra}). Por lo tanto, ella o un punto que la domine (con $f_3$
al menos igual) queda en el frente devuelto. En ambos casos,
$\max_{\Front}f_3=b^\ast/n$, y todo punto de $\Front_{\mathrm{s}}$ tiene precio
en las $b^\ast$ líneas que lo admiten. Las cuatro proyecciones eligen dentro de
$\Front_{\mathrm{s}}$.
\end{proof}

% ---------------------------------------------------------------------------
\subsection{De la asignación al kit: líneas, herramientas y consolidación}
\label{sec:opt-lineas}
% ---------------------------------------------------------------------------

Elegida la asignación, el kit se construye con tres transformaciones que no
dependen del algoritmo que eligió.

\paragraph{1. Construcción de líneas.} Cada renglón informativo produce una
línea sin precio ni subtotal, con la instrucción del colegio. La base de datos
impone que esas líneas no tengan precio. Cada renglón decidible con candidato
$c$ produce una línea con $u_i(c)$ unidades, subtotal $\kappa_i(c)$ (o
$\sinprecio$) y hasta dos alternativas. Un renglón decidible sin candidato
produce una línea vacía con su mensaje.

\paragraph{2. Juntar herramientas de un mismo hijo.} Las listas se escriben por
materia, y cada profesor pide lo que usa: las tijeras aparecen en artes, en
tecnología y en la cartuchera. Son las mismas tijeras. En las listas curadas,
43 renglones de herramientas repetían a otro idéntico; un grado pedía tres
calculadoras científicas. Para los conceptos marcados en el catálogo como
\emph{uno por estudiante}, las líneas de un mismo hijo con el mismo producto
elegido se agrupan:
\[
  \text{grupo } g=\{i : h(r_i)=h,\ x_i=c,\ c \text{ es uno por estudiante}\},
  \qquad q_g=\max_{i\in g}q(r_i).
\]
Se compra la cantidad \emph{máxima}, no la suma: si artes pide dos reglas y la
cartuchera una, dos alcanzan. La primera línea del grupo lleva
$\techo{q_g/\emp(c)}$ unidades y las demás quedan con cero, con un mensaje que
dice en qué línea quedó el producto. La agrupación es por \emph{hijo}, porque
dos hermanos necesitan cada uno sus tijeras, y por \emph{producto elegido},
porque «regla» y «regla de acero» no son la misma compra. Se limita a los
conceptos uno por estudiante porque lo que se gasta (pegante, cinta,
cuadernos) puede pedirse uno por materia a propósito, y juntarlo dejaría al
niño sin material.

\paragraph{3. Consolidar entre hermanos.} Sobre las líneas resultantes, las de
\emph{distintos} hijos con el mismo producto elegido se consideran para una
compra conjunta, excluyendo los textos (dos ediciones del mismo título no son
sustituibles) y las líneas ya juntadas en el paso 2. Para un grupo $g$ con
empaque $\emp$, se comparan
\[
  \text{juntas}=\techo{\frac{\sum_{i\in g}q(r_i)}{\emp}}
  \qquad\text{y}\qquad
  \text{sueltas}=\sum_{i\in g}\techo{\frac{q(r_i)}{\emp}},
\]
y se consolida solo si $\text{juntas}<\text{sueltas}$. En ese caso, la primera
línea lleva las unidades conjuntas, las demás quedan en cero con su
explicación, y se registra un ahorro de $\precio\cdot(\text{sueltas}-\text{juntas})$.

\begin{lema}[Subaditividad del techo]
\label{lem:techo}
Para $a,b\in\N$ y $\emp\in\N_{\geq1}$,
$\techo{(a+b)/\emp}\leq\techo{a/\emp}+\techo{b/\emp}$.
\end{lema}

\begin{proof}
$\techo{a/\emp}+\techo{b/\emp}$ es un entero mayor o igual que
$a/\emp+b/\emp=(a+b)/\emp$, y $\techo{(a+b)/\emp}$ es el menor entero con esa
propiedad.
\end{proof}

\begin{corolario}[Consolidar nunca encarece]
\label{cor:consolidar}
Para cualquier grupo, $\text{juntas}\leq\text{sueltas}$ (por inducción sobre
el \cref{lem:techo}), así que la consolidación nunca aumenta el total del kit.
La condición estricta $\text{juntas}<\text{sueltas}$ evita además mover
unidades entre hermanos cuando no hay beneficio, lo que solo haría más difícil
leer la explicación.
\end{corolario}

\begin{ejemplo}
Dos hermanos piden un borrador cada uno y el producto elegido es una caja de
10 a \$9.000: sueltas $=1+1=2$ cajas, juntas $=\techo{2/10}=1$ caja, ahorro
$=\$9.000$.
\end{ejemplo}

La consolidación se hace sobre el \emph{producto} elegido y no sobre el
concepto. Si a dos hermanos se les eligieron productos distintos, forzarlos a
compartir empaque cambiaría lo que recibe cada uno sin decirlo. Las líneas que
quedan en cero no se borran: el colegio pidió ese ítem para ese hijo, y el
acudiente tiene que poder encontrarlo. Esta es la razón de fondo para que una
solicitud sea por \emph{hogar}: si cada hijo se pidiera por separado, nadie
vería que una caja de diez sale mejor que dos sueltos.

El \emph{total} del kit es la suma de los subtotales de las líneas no
informativas. Lo que el sistema no puede entregar nunca se cobra.

\begin{observacion}[Lo que el optimizador no ve]
Las dos agrupaciones se aplican \emph{después} de elegir, así que $f_1$ no las
ve. Una asignación que dé el mismo producto a dos hermanos puede resultar más
barata, después de consolidar, que otra que el frente consideraba
equivalente. Este es el acoplamiento que el diseño original proponía llevar
dentro del objetivo (\cref{sec:opt-estructura}).
\end{observacion}

% ---------------------------------------------------------------------------
\subsection{Qué kits se arman y el kit «A tu Medida»}
\label{sec:opt-medida}
% ---------------------------------------------------------------------------

No se arman todos los kits posibles, sino los que el acudiente pidió:
\begin{itemize}
  \item \textbf{Ahorrador} y \textbf{Premium}, siempre.
  \item \textbf{Fan}, solo si se pidió un personaje para algún hijo.
  \item \textbf{A tu Medida}, solo si se eligió la prioridad de
        \emph{equilibrio} o se declaró un tope de gasto.
\end{itemize}
Un kit que nadie pidió es ruido en la pantalla y trabajo de más para el
revisor. Los kits \emph{canónicos} de una lista, que se arman sin
preferencias para reutilizarlos, son por eso solo dos.

\begin{definicion}[A tu Medida]
\label{def:a-tu-medida}
Sea $T$ el tope: el presupuesto declarado $B$, si existe; si no,
$T=\piso{(\mathrm{total}(k_{\mathrm{Ah}})+\mathrm{total}(k_{\mathrm{Pr}}))/2}$, la
mitad del camino entre el Ahorrador y el Premium. El kit es
\[
  \mathbf{x}_{\mathrm{AM}}=\lexmax_{\mathbf{x}\in\Front_{\mathrm{s}},\ f_1(\mathbf{x})\leq T}\ \bigl(f_2,\ -f_1,\ f_3\bigr),
\]
con $\Front_{\mathrm{s}}$ los puntos surtibles de la \cref{def:proyeccion},
y no se arma (se devuelve $\varnothing$) si
\begin{enumerate*}[label=(\roman*)]
  \item el frente es vacío,
  \item ningún punto cabe,
  \item el total del kit construido, ya consolidado, supera $T$, o
  \item el kit tiene exactamente los mismos productos que el Ahorrador o que
        el Premium.
\end{enumerate*}
\end{definicion}

Es la mejor marca que cabe en el tope, tomada del frente ya calculado. La
condición (iv) evita mostrar el mismo kit dos veces con nombres distintos. El
presupuesto se aplica \emph{después} del frente y no como restricción del
NSGA-II: un frente de Pareto no debe esconder un kit válido porque es caro,
sino ofrecerlo para que la capa de recomendación elija y lo explique. Meter el
presupuesto en la optimización recortaría el frente. El tope es el presupuesto
tal cual: la disposición a pagar sobreprecio no lo estira, porque indica
cuánto más pagaría la familia por marca \emph{sin} tope, no cuánto se puede
pasar del tope que ella misma puso.

% ---------------------------------------------------------------------------
\subsection{La recomendación}
\label{sec:opt-recomendar}
% ---------------------------------------------------------------------------

Entre los kits armados se marca \emph{uno} como recomendado, sin ocultar los
demás. La regla es una lista de decisión cuyas condiciones salen de una
encuesta a 60 familias:

\begin{enumerate}
  \item Si se pidió un personaje y existe el Fan $\Rightarrow$ \textbf{Fan}
        («es el único kit que lo trae»).
  \item Si existe A~tu~Medida $\Rightarrow$ \textbf{A tu Medida}, con una
        razón que dice si el tope lo puso el acudiente o es la mitad del camino.
  \item Si se pidió el punto medio pero A~tu~Medida no se armó
        $\Rightarrow$ \textbf{Premium}, si su total cabe en el tope, o
        \textbf{Ahorrador} en otro caso, con la razón correspondiente.
  \item Si el peso declarado para la calidad es al menos el de cada otro
        criterio y la disposición a pagar sobreprecio es $\sigma\geq0{,}10$
        $\Rightarrow$ \textbf{Premium}.
  \item En otro caso $\Rightarrow$ \textbf{Ahorrador}. Es el valor por
        omisión porque 42 de las 60 familias encuestadas pagarían como
        máximo un 5\,\% más, y 15 exigen un precio igual o menor.
\end{enumerate}

\paragraph{El presupuesto como restricción dura.} Si hay un tope $B$ y el kit
marcado no cabe, la marca se mueve al kit \emph{de mejor marca que sí cabe}:
\[
  k^\star=\lexmax_{k:\ \mathrm{total}(k)\leq B}\ \bigl(f_2(k),\ -\mathrm{total}(k)\bigr).
\]
Si ninguno cabe, se marca el más barato y la razón dice cuánto cuesta. Antes
la marca caía al kit más barato, y eso dejaba al recomendado con presupuesto en
el 40\,\% de la calidad posible.

\paragraph{Relación con la teoría de la utilidad multiatributo.} El diseño
original preveía que la recomendación fuera una escalarización aditiva
$U(k)=\sum_jw_ju_j(k)$ de utilidades normalizadas sobre el rango observado en
el frente, en la línea de la teoría de la utilidad multiatributo
\cite{keeney1976}, con las reglas de la encuesta como distribuciones previas
sobre los pesos. La versión vigente implementa las reglas directamente como
lista de decisión, y los pesos declarados solo intervienen en la condición
«calidad primero» del paso 4. La lista de decisión es más fácil de auditar y
de explicar al acudiente. La escalarización permitiría usar toda la
información de los pesos, pero exige suponer independencia preferencial
entre los criterios, que es lo que justifica la forma aditiva. Hay además una
diferencia documentada entre los pesos por omisión del código
(precio 0,35; marca exigida 0,30; calidad 0,25; personaje 0,10) y los que
resultan de normalizar la pregunta de prioridades de la encuesta
(0,285; 0,273; 0,255; 0,187). Los del código exageran la separación entre
criterios y le quitan al personaje casi la mitad de su peso medido.

% ---------------------------------------------------------------------------
\subsection{El personaje: inyección y anclaje}
\label{sec:opt-personaje}
% ---------------------------------------------------------------------------

Los productos con personaje no compiten en la resolución
(\cref{sec:res-niveles}), así que el frente no siempre tiene un candidato con
personaje para elegir. Cuando el acudiente pide un personaje, dos
post-procesos sobre el kit Fan cierran ese hueco:

\paragraph{Inyección.} Para cada concepto \emph{ancla} (cuaderno, cartuchera,
carpeta y maleta) que el catálogo surte con el personaje pedido, el producto
elegido en la línea del Fan se reemplaza por el producto con personaje. Como
el reemplazo cambia precios después de que el kit se armó, cada línea se
devuelve a su forma previa a la consolidación, se vuelve a consolidar ---dos
hermanos que comparten un ancla deben volver a compartir empaque--- y se
recalculan el total, los objetivos y la recomendación completa. De lo
contrario, el kit mostraría un total nuevo junto a objetivos y un veredicto
de presupuesto calculados con los precios viejos.

\paragraph{Anclaje.} Las líneas ancla cuyo producto trae el personaje pedido
se mueven al principio del kit, en el orden canónico cuaderno, cartuchera,
carpeta, maleta. El resto conserva su orden relativo; es una partición
estable.

\begin{proposicion}[El anclaje no cambia el kit]
El anclaje preserva el total, las consolidaciones y los objetivos.
\end{proposicion}

\begin{proof}
El anclaje es una permutación de las líneas. El total y los objetivos son
sumas y promedios sobre las líneas, invariantes bajo permutaciones, y las
consolidaciones ya se calcularon y no se recalculan.
\end{proof}

% ---------------------------------------------------------------------------
\subsection{El algoritmo completo y su complejidad}
\label{sec:opt-completo}
% ---------------------------------------------------------------------------

\begin{algorithm}[htbp]
\caption{Armar los kits de un hogar}
\label{alg:armar}
\begin{algorithmic}[1]
\Require renglones $R$ con candidatos; preferencias; semilla $\varsigma$
\Ensure kits armados, uno marcado como recomendado
\State $\Dec\gets$ renglones no informativos con $A(r)\neq\varnothing$;\quad $\Esp\gets\prod_iA_i$
\If{$\Dec=\varnothing$ \OR $\card{\Esp}\leq1$} $\Front\gets\varnothing$
\ElsIf{$\card{\Esp}\leq4096$} $\Front\gets\textsc{FrenteExacto}(\Dec)$ \Comment{\cref{alg:frente-exacto}}
\Else
  \State $S\gets$ heurística por perfil $\cup$ \textsc{Escalera}$(\Dec)$ \Comment{\cref{def:heuristica,alg:escalera}}
  \State $\Front\gets\mathrm{NoDom}\bigl(\textsc{NSGA-II}(N(n),G(\Esp),S,\varsigma)\cup S\bigr)$ \Comment{\cref{alg:nsga2}}
\EndIf
\For{perfil $\in\{\text{Ahorrador},\text{Premium}\}$ y, si se piden, Fan}
  \State $\mathbf{x}\gets$ proyección de $\Front$ (\cref{def:proyeccion}), o heurística si $\Front=\varnothing$
  \State $k_{\text{perfil}}\gets\textsc{Consolidar}(\textsc{JuntarHerramientas}(\textsc{Líneas}(R,\mathbf{x})))$
\EndFor
\If{se pidió el punto medio o un tope} armar A~tu~Medida (\cref{def:a-tu-medida}) \EndIf
\State marcar el recomendado (\cref{sec:opt-recomendar})
\If{se pidió un personaje} inyectar y anclar en el Fan; volver a recomendar \EndIf
\State \Return los kits
\end{algorithmic}
\end{algorithm}

\paragraph{Complejidad global.} Sea $N_R=\card{R}$. La construcción de líneas,
la agrupación de herramientas y la consolidación cuestan $O(N_R)$ por kit con
tablas de dispersión, y se arman como mucho cuatro kits. La proyección es
lineal en $\card{\Front}$. El costo lo domina la construcción del frente:
\[
  T_{\mathrm{opt}}=\begin{cases}
    O(nk) & \card{\Esp}\leq1,\\
    O(\card{\Esp}\,n+\card{\Esp}^2) & \card{\Esp}\leq4096 \quad(\leq 1{,}7\times10^7),\\
    O\bigl(G(MN^2+Nn)+n^2k\bigr) & \text{en otro caso.}
  \end{cases}
\]
En la evaluación sobre 29 listas, producir todos los kits de una lista tomó en
promedio 5,3~s, con mediana de 6,3~s y máximo de 9,0~s, en un solo proceso de
CPU. Es el mismo orden de magnitud que leer una página con el modelo
multimodal.

% ---------------------------------------------------------------------------
\subsection{Evolución del optimizador}
\label{sec:opt-evolucion}
% ---------------------------------------------------------------------------

\begin{evolucion}[: optimizador, versión por versión]
\begin{description}[style=nextline,leftmargin=0pt]
  \item[v0 — Prototipo heredado.] Un NSGA-II con cuatro objetivos que eran
        sumas sobre ítems independientes, una variable por ítem y
        $\text{costo}=\text{precio}\times\text{cantidad}$. El problema era
        completamente separable, y sus propias pruebas lo confirmaban: varias
        afirmaban que el NSGA-II devolvía \emph{exactamente} lo mismo que la
        fuerza bruta, lo que indicaba que no había nada que buscar. Tenía
        además tres defectos: no había empaque (el error de las doce cajas),
        con presupuesto se duplicaba el kit más barato hasta completar tres
        (el acudiente podía recibir tres kits idénticos presentados como tres
        opciones) y uno de sus objetivos valía siempre 1.
  \item[v1 — Andamio heurístico.] Una regla de selección por perfil,
        determinista y declarada como provisional en su propia documentación.
        Su propósito era que el ciclo completo corriera y se pudiera probar
        antes de que cada etapa fuera óptima. De esta versión se conservan sin
        cambios la aritmética de unidades, las líneas informativas que no se
        cobran y la forma de la recomendación.
  \item[v2 — Primer NSGA-II.] Frente de Pareto con NSGA-II y enumeración
        exacta para $\card{\Esp}\leq4096$. Las variables eran \emph{reales} en
        $[0,k_i-10^{-6}]$ con $x_i\mapsto\piso{x_i}$, la relajación continua
        estándar, con población fija de 60 y 40 generaciones.
  \item[v3 — Codificación entera.] Las variables pasaron a ser enteras en
        $[0,k_i-1]$ con reparación por redondeo, porque con la relajación
        continua SBX interpolaba un índice nominal como si fuera una magnitud.
        La población y las generaciones pasaron a escalarse con el problema.
        La biblioteca de optimización pasó a ser una dependencia obligatoria
        (antes se degradaba en silencio a la heurística). La «fidelidad» que se
        mostraba pasó a ser el mismo eje que se optimiza; antes era la
        fracción de líneas resueltas automáticamente, y los dos números solo
        coincidían por correlación.
  \item[v4 — Correcciones de la evaluación.] Tres cambios motivados por las
        mediciones de la \cref{sec:evaluacion}:
        \begin{enumerate*}[label=(\alph*)]
          \item la heurística ordena por costo de línea y no por precio
                unitario (\cref{lem:precio-unitario});
          \item el NSGA-II se siembra con la escalera de costo marginal, que
                también entra al filtro final (\cref{prop:siembra}); y
          \item con presupuesto, se recomienda la mejor marca que cabe,
                tomada del frente.
        \end{enumerate*}
        Efecto medido en las 174 listas (semilla 0): la heurística pasó de
        acertar el 80\,\% de los Ahorrador al 100\,\%, la escalera pasó de
        dominar algún punto del frente en el 68\,\% de las listas al 0\,\%, y
        el recomendado con presupuesto pasó del 40\,\% de la calidad posible a
        caber siempre e igualar el óptimo exacto.
  \item[v5 — Kits a pedido.] Ahorrador y Premium siempre; Fan y A~tu~Medida
        solo si se piden. El kit con tope dejó de llamarse «Premium ajustado» y
        pasó a ser A~tu~Medida (\cref{def:a-tu-medida}). El tope dejó de
        estirarse con la disposición a pagar sobreprecio: un acudiente que
        escribía «máximo \$350.000» podía recibir un kit de \$385.000.
  \item[v6 — Herramientas por hijo.] Agrupación de las herramientas pedidas
        por varias materias (\cref{sec:opt-lineas}), motivada por los 43
        renglones repetidos en las listas curadas.
  \item[v7 — Proyección sobre lo surtible.] Al formalizar la proyección para
        este documento apareció que un candidato sin precio parecía gratis en
        el frente: el Ahorrador podía elegirlo con un total que no incluía esa
        línea, y el Premium y A~tu~Medida podían tomarlo como una marca sin
        costo. Desde entonces, las cuatro proyecciones eligen entre los puntos
        de máxima disponibilidad (\cref{def:proyeccion,prop:surtibles}). El
        catálogo vigente no tenía productos activos sin precio, así que ningún
        kit cambió; la corrección cierra un defecto latente que habría
        aparecido con la primera oferta vencida.
\end{description}
\end{evolucion}
